\subsection{Test harness}
\label{sec:method-harness}

Every server is built as a Docker image pinning the exact version
under test and exposing port 8000 on a Docker bridge network. The
container is constrained to a single CPU
(\texttt{--cpus=1}). The application handler implements the
minimum needed for a meaningful measurement: a \texttt{GET /health}
returning \texttt{200 OK} with body \texttt{\{"ok":true\}} and a
\texttt{POST /upload} that fully drains the request body and
returns \texttt{200 OK} with body \texttt{\{"len":N\}}. No logging
middleware, no metrics, no compression, no body-size limit is
configured beyond the framework default.

A separate \emph{prober} container, on the same Docker bridge
network, reaches the server by its DNS alias and issues the
controlled HTTP requests. The prober is a single Python script that
opens a raw \texttt{asyncio} TCP connection, writes the request
headers, then writes the body as chunked-TE chunks per the test
mode (see below), and reads the response. Wall time is measured
client-side from the moment the prober writes the first byte of
the body to the moment the server's response is received.

\subsection{Standardized payload}
\label{sec:method-payload}

For each measurement cell we send a body of $N$ bytes, encoded as
$N$ chunks of one byte each:
\begin{verbatim}
POST /upload HTTP/1.1
Host: x
Transfer-Encoding: chunked
Content-Type: application/octet-stream
Connection: close

1\r\nA\r\n1\r\nA\r\n... (N times) ...0\r\n\r\n
\end{verbatim}
Each chunk is \texttt{1$\backslash$r$\backslash$nA$\backslash$r$\backslash$n}
(6 wire bytes per data byte), so the bytes on the wire are roughly
$6N$ plus a fixed header overhead. We measure at
$N \in \{50000, 100000, 250000\}$, with a few servers extended to
$N{=}1000000$ for the log-log scaling charts in
Figure~\ref{fig:scaling}.

This payload is RFC-compliant. It is a degenerate-but-legal
chunked-TE encoding; no protocol-level violation is required.

\subsection{Two probe modes}
\label{sec:method-modes}

The TCP coalescing behavior on the server depends on how the prober
delivers bytes to the kernel. We measure both extremes:

\paragraph{Mode A: bridge-coalesced.}
The prober writes all chunks back-to-back to its
\texttt{StreamWriter}; \texttt{TCP\_NODELAY} is \emph{not} enabled.
The kernel coalesces the writes into MSS-sized TCP segments. The
server's \texttt{recv()} buffer receives many chunks per syscall
return. Empirically, on our Docker bridge this is $\sim$8--16
wire chunks per \texttt{data\_received} callback.

\paragraph{Mode B: paced \SI{100}{\micro\second}.}
Between every chunk write the prober busy-waits for
\SI{100}{\micro\second} (a CPU spin, not \texttt{asyncio.sleep};
the latter yields to the event loop and is too coarse).
\texttt{TCP\_NODELAY} is enabled. Each chunk leaves the prober as
its own TCP segment with $\geq$\SI{100}{\micro\second} between
sends; the server typically returns from each \texttt{recv()} with
one chunk's worth of bytes.

Mode A and Mode B bound the realistic spectrum of network
conditions: Mode A approximates kubernetes pod-to-pod traffic,
Mode B approximates slow-drip attacker pacing or trans-WAN bursts.

\subsection{What the per-chunk cost number means}
\label{sec:method-metric}

\paragraph{Primary observable.}
The one quantity we record on every cell is \emph{prober-side wall
time} $T_\mathrm{wall}$: elapsed seconds from the prober writing
the first body byte to receiving the server's response. This is an
end-to-end elapsed-time measurement, not a CPU-time measurement.

\paragraph{Per-chunk cost, defined per mode.}
For a cell of $N$ one-byte chunks we report a per-chunk cost $c$ in
microseconds. Because the two modes load the server differently,
$T_\mathrm{wall}$ contains different components, so we define $c$
per mode:
\begin{itemize}
  \item \textbf{Mode A (bridge-coalesced):} no inter-chunk delay,
    so $T_\mathrm{wall}$ is dominated by server work plus
    scheduling. We report $c_A = T_\mathrm{wall}\cdot 10^{6}/N$, a
    wall-derived upper bound on per-chunk server cost (it includes
    any prober-side and network time, and for the quadratic
    servers it is the steady-state slope, not a constant).
  \item \textbf{Mode B (paced \SI{100}{\micro\second}):} the
    prober inserts a \SI{100}{\micro\second} busy-wait per chunk,
    so the floor of $T_\mathrm{wall}$ is
    $N\cdot\SI{100}{\micro\second}$ whenever the server is faster
    than the pacing. The per-chunk \emph{server} cost is then the
    excess over that floor,
    $c_B = (T_\mathrm{wall}-N\cdot\SI{100}{\micro\second})\cdot
    10^{6}/N$. When the server's own per-chunk cost \emph{exceeds}
    \SI{100}{\micro\second} the pacing is fully overlapped and no
    longer the floor; for those saturated cells
    ($\mathrm{cpu}\!\approx\!100\%$, e.g.\ \fwid{nginx},
    \fwid{bandit}) we instead report
    $T_\mathrm{wall}\cdot 10^{6}/N$, a wall-time upper bound rather
    than a clean CPU figure.
\end{itemize}

\paragraph{Wall-derived vs.\ CPU-derived cells.}
Where the container is saturated for the full run
($\mathrm{cpu}\!\approx\!100\%$), wall and CPU time coincide to
within measurement error and we treat the figure as CPU.
Otherwise wall time over-counts CPU, so the per-chunk number is a
wall-time upper bound. Of the 203 cells in our dataset, only 64
carry a server CPU sample (\texttt{server\_cpu\_pct}); the other
139 are wall-derived. We never convert a wall-derived cell into a
CPU-second claim, and Section~\ref{sec:threats} records the
resulting bias directions.

\paragraph{CPU sampling method.}
Where present, \texttt{server\_cpu\_pct} is a single
\fname{docker stats --no-stream} snapshot taken during the run,
and the CPU-seconds we cite are
$T_\mathrm{wall}\times\texttt{server\_cpu\_pct}$. This is a
sampled rate times a duration, not an integral of
\fname{cpuacct.usage}; under contention \fname{docker stats} can
over- or under-report and can exceed \SI{100}{\percent} on a
single pinned core, so we round CPU figures to the
tens-of-percent level.

\subsection{Profile collection}
\label{sec:method-profile}

For each server we sample its CPU profile during the worst Mode B
cell, using whichever profiler the language ecosystem provides:

\begin{itemize}
  \item Python: \texttt{cProfile} with \texttt{sort\_stats('cumulative')},
    top 30.
  \item Go: \texttt{net/http/pprof} exposed on a side port, sampled
    via \texttt{go tool pprof -top}.
  \item Rust: \texttt{perf record -F 999}, converted to
    \texttt{perf report --stdio} (\texttt{kernel.perf\_event\_paranoid}
    blocked profiling on some Rust runs; we fall back to
    source-level analysis for those).
  \item Node: \texttt{node --prof}, then
    \texttt{node --prof-process} top 100.
  \item JVM: \texttt{async-profiler} agent attached, itimer
    sampling.
  \item BEAM: \texttt{:eprof} posted to a \texttt{/profile\_upload}
    endpoint inside the same process.
  \item .NET: \texttt{dotnet-trace} via a sidecar container
    sharing the target's PID namespace.
\end{itemize}

\subsection{Reproducibility}
\label{sec:method-repro}

Every test harness is checked into the project tree under
\texttt{wave\{1,2\}/<ecosystem>/repro/}, including the
Dockerfile(s), the application source, the prober Python script,
the per-cell run logs, and the saved profile artifacts. The
top-level rebuild driver,
\fname{scripts/render\_all.sh}, regenerates every figure from
\texttt{data/wave*.json} via the chart-rendering scripts and
re-runs the LuaLaTeX build of this paper. Anyone with Docker plus
\fname{uv} plus \fname{lualatex} can reproduce every number in
this paper from the published harness.
